\subsection{Empirical implications}\label{app:empirical}

This subsection states and proves the results that Section~\ref{sec:empirical} summarizes. The auction is the benchmark, the prior is
\(1/2\), and the slot lies in the band of Proposition~\ref{prop:region}(a). The price increment of entry
\(A_r(\mu)=t_L-t_0+\DT(r)\mu\) is the function \(w(\mu)\) of Appendix~\ref{app:beliefs} at the strength \(r\). It is
positive, rises in \(\mu\) with slope \(\DT\), and rises in \(r\), because \(t_L-t_0=\tfrac12-1/(8r)\) and \(\DT\) rise
in \(r\) at the benchmark. The indicator \(\chi_r=\1\{B_r(M)\ge c_H\}\) records whether a purchase recruits, and the
volume \(v=\Pr(X\ne0)\) is the frequency of nonzero orders per order opportunity. Write \(\bar f=b+\lambda/2\), \(\vartheta=u/(2+u)=2M-1=M-m\), and \(e_+=\rho+(1-\rho)\chi_r\). At the
benchmark, \(B_r(1/2)=5-r/4+1/(8r)<c_H\) for \(1<r<h\). The top profit \(B_r(M)\) falls strictly in \(r\) for
\(r>\ell\), equals \(9M+3/8\) at \(r=\ell\), and lies below \(c_H\) at \(r=h\); so for \(M>5/8\) the ceiling \(r_C\) is the
unique strength at which \(B_r(M)=c_H\), and \(\chi_r=1\) exactly when \(r\le r_C\). In the band, \(M-5/8=\reg{margin_obs_post_over_58_B}\)
at the benchmark slot.

\paragraph{Full two-sided trade on \((r_S,h)\).}
The strength \(r_S\) solves \(\rho\DT(r)m=k\) exactly: \(r_S=\fr(D_S)\) with \(D_S=k/(\rho m)=\reg{delta_rS_B}\). For
\(r_S<r<h\), every candidate satisfies
\[
K_H\ge\rho\DT(r)(1-M)=\rho\DT(r)m>k,\qquad K_L\ge\rho\DT(r)m>k
\]
by \eqref{eq:a2-bounds}, since \(1-M=m\) at the prior \(1/2\) and \(\DT\) rises. So full orders are the unique
equilibrium orders, under both investor rules. The tuple \(\mathcal F\) has the posteriors \((m,1/2,M)\), entry
\((\rho,\rho,e_+)\), the prices \(P(-1)=t_0+\rho A_r(m)<P(0)=t_0+\rho A_r(1/2)<P(1)=t_0+e_+A_r(M)\), \(\bar\mu(P(x))=\mu(x)\),
\(\zeta(P(1))=\chi_r\), and truthful bids. Each condition holds by direct computation. (E1): the likelihoods
\(f_H=(b,b_0,b+\lambda)\) and \(f_L=(b+\lambda,b_0,b)\) give \((m,1/2,M)\). (E2): entry is nondecreasing in the posterior,
so the three prices differ, and each price reveals its order. (E3): \(B_r(m)<B_r(1/2)<c_H\), and the purchase
recruits exactly when \(\chi_r=1\). (E4): \eqref{eq:a2-terminal}. (E5): the strict bounds above. (E6): truthful bids. At \(r_C\) rule C1 gives \(\chi_{r_C}=1\). The interval
contains \(r_1\), \(r_C\) and \(r_2\), because \(r_S<r_1<r_C<r_2<h\) by Lemma~\ref{lem:a4-order}; at the benchmark
\(r_C-r_S=\reg{margin_obs_rS_rC_B}\). The bound is exact: \(r_S\approx\reg{rS_B}\), and the six-place display lies below
\(r_S\). At \(r=\reg{rS_counter_r_B}\), inside the rounded interval, \(\rho\DT m-k=\reg{rS_counter_margin_B}\), so the low
class strictly prefers abstention to the full sale there; \(\reg{rS_cert_hi_B}\) is a sufficient rounded endpoint. With
a hidden price, entry is \(\rho\) after every order under every profile, because the belief stays at \(1/2\) and
\(B_r(1/2)<c_H\); the same bounds with \(e\equiv\rho\) force full orders, and the tuple is \(\mathcal F\) with
\(\chi_r\) replaced by \(0\).

Under full orders, class \(H\) places the order \(1\) with probability \(b+\lambda\) and class \(L\) with probability
\(b\), so \(\Pr(X=1)=\Pr(X=-1)=\bar f\), \(\Pr(X=0)=1-2\bar f\), \(v=2b+\lambda\) and \(\E[X]=0\). The informed share of nonzero
orders is \(\lambda/(2b+\lambda)=u/(2+u)=\vartheta\). Only the order \(1\) can recruit, so entry is
\(\rho+(1-\rho)\chi_r\bar f\); from \(r_0\), where the orders are zero and entry is \(\rho\), to \(r_1\) it rises by
\((1-\rho)\bar f=\reg{entry_jump_B}\).

The spread \(S=P(1)-P(-1)\) is the difference between the conditional execution prices of a unit purchase and a
unit sale. The purchase impact is \(I_+=P(1)-P(0)\), and the sale impact is \(I_-=P(0)-P(-1)\). The model has no
separate quote stage: each price is the competitive value conditional on the order.

\begin{proposition}[The spread and its entry component; analytical]\label{prop:spread}
Let \(r_S<r<h\).
\begin{enumerate}[label=(\roman*),leftmargin=2em]
\item \(S=\rho\DT\vartheta+(1-\rho)\chi_rA_r(M)\): a class component and an entry component.
\item \(I_-=\rho\DT\vartheta/2\) and \(I_+=I_-+(1-\rho)\chi_rA_r(M)\). Hence \(I_+>I_-\) if and only if a purchase
recruits.
\item If the challenger does not see the price, full two-sided trade is still the unique outcome, and \(S\) equals
the class component \(\rho\DT\vartheta\).
\item On \((r_S,r_C]\), both components, both impacts and \(S\) rise strictly with \(r\). On \((r_C,h)\), the entry
component is zero, and the class component, both impacts and \(S\) rise strictly.
\item At \(r_C\), \(S\) falls by \((1-\rho)A_{r_C}(M)>0\). The prices \(P(-1)\) and \(P(0)\) and the sale impact are
continuous there; \(P(1)\) and the purchase impact fall by this amount.
\end{enumerate}
\end{proposition}

\begin{proof}[Proof of Proposition~\ref{prop:spread}]
Let \(r_S<r<h\), so the outcome is \(\mathcal F\).
\begin{enumerate}[label=(\roman*),leftmargin=2em]
\item Subtract the prices: \(S=\rho[A_r(M)-A_r(m)]+(1-\rho)\chi_rA_r(M)\). Since \(A_r\) has slope \(\DT\),
\(A_r(M)-A_r(m)=\DT(M-m)=\DT\vartheta\).
\item Also \(A_r(1/2)-A_r(m)=\DT(1/2-m)=\DT\vartheta/2\) and \(A_r(M)-A_r(1/2)=\DT\vartheta/2\), so
\(I_-=\rho\DT\vartheta/2\) and \(I_+=\rho\DT\vartheta/2+(1-\rho)\chi_rA_r(M)\). Since \(A_r(M)>0\), \(I_+>I_-\) exactly when
\(\chi_r=1\).
\item With a hidden price the outcome is \(\mathcal F\) with entry \(\rho\) at every order, so part (i) holds with
\(\chi_r\) replaced by \(0\).
\item On \((r_S,r_C]\), \(\chi_r=1\); \(\DT\) and \(A_r(M)\) rise strictly in \(r\), and \(\rho\), \(\vartheta\) and
\(1-\rho\) do not depend on \(r\). So both components, both impacts and \(S\) rise strictly. On \((r_C,h)\), \(\chi_r=0\):
the entry component is identically zero, and the class component \(\rho\DT\vartheta\), the impacts
\(I_-=I_+=\rho\DT\vartheta/2\) and \(S=\rho\DT\vartheta\) rise strictly.
\item Since \(r_S<r_C\), both sides of \(r_C\) lie in \(\mathcal F\). Rule C1 gives \(\chi_{r_C}=1\), and \(\chi_r=0\) for
\(r>r_C\). Every other term is continuous in \(r\). So \(P(-1)\), \(P(0)\) and \(I_-\) are continuous at \(r_C\), and
\(P(1)\), \(I_+\) and \(S\) fall by \((1-\rho)A_{r_C}(M)>0\).\qedhere
\end{enumerate}
\end{proof}
In the benchmark at \(r_1\) the class component is \(\reg{ba_class_r1_B}\) and the entry component \(\reg{ba_entry_r1_B}\),
so \(S=\reg{ba_r1_B}\); at \(r_C\) the spread falls from \(\reg{ba_rC_left_B}\) to \(\reg{ba_rC_right_B}\). Under rule CM,
entry after a purchase at \(r_C\) can take any value in \([\rho,1]\), and \(S(r_C)\) can take any value between its
two one-sided limits. The spread is the difference between conditional execution prices; the model has no
separate quote stage.

\begin{proposition}[Identification of the informed share; analytical]\label{prop:informedshare}
Let \(r_S<r<h\), and write \(v=2b+\lambda\).
\begin{enumerate}[label=(\alph*),leftmargin=2em]
\item \emph{Counts.} The order law depends on \((\lambda,\beta)\) only through \(v\), and entry
\(\rho+(1-\rho)\chi_rv/2\) depends on them only through \(v\) and \(\chi_r\). For a ratio \(u'>0\), a slot
\((\lambda',\beta')\) with this ratio and the volume \(v\) exists if and only if \(3v\le2+u'(1-v)\). Equivalently, its informed
share \(\vartheta'=u'/(2+u')\) satisfies \(\vartheta'\ge3-2/v\). This slot is then unique.
\item \emph{Incentives.} A slot with informed share \(\vartheta'\) produces these counts through full two-sided trade
at \(r\) only if \(\vartheta'\le\bar\vartheta(r)=1-2k/(\rho\DT(r))\).
\item \emph{Band.} The slots with strict recruitment at \(r_1\) and strict exclusion at \(r_2\) have informed shares
exactly in \((2\tau(r_1)-1,\,2\tau(r_2)-1)\). At the benchmark
volume, each share in this band belongs to a slot with the same counts and the same entry at \(r_1\) and at \(r_2\).
Full two-sided trade is the unique outcome of this slot at both strengths.
\item \emph{Prices.} Given the entry floor \(\rho=\Pr(\text{preparation}\mid X\ne1)\) and the spread \(\DT(r)\), the
sale impact and the volume identify the slot: \(\vartheta=2I_-/(\rho\DT(r))\), \(\lambda=v\vartheta\) and
\(\beta=v(1-\vartheta)/[2(1-v\vartheta)]\).
\end{enumerate}
\end{proposition}

\begin{proof}[Proof of Proposition~\ref{prop:informedshare}]
Let \(r_S<r<h\).
\begin{enumerate}[label=(\alph*),leftmargin=2em]
\item The order law of \(\mathcal F\) is \((v/2,1-v,v/2)\) with \(v=2b+\lambda\), and entry is \(\rho+(1-\rho)\chi_rv/2\).
Given \(v\) and \(u'>0\), the equations \(v=\lambda'+2b'\) and \(u'=\lambda'/b'\) give \(b'=v/(u'+2)\),
\(\lambda'=vu'/(u'+2)\) and \(\beta'=b'/(1-\lambda')=v/[2+u'(1-v)]\), a unique pair. It has \(0<\lambda'<v<1\) and
\(\beta'>0\), so it is a slot of the model exactly when \(\beta'\le1/3\), that is \(3v\le2+u'(1-v)\). With
\(u'=2\vartheta'/(1-\vartheta')\), the right side is \((2-2v\vartheta')/(1-\vartheta')\), and the condition reads
\(3v(1-\vartheta')\le2-2v\vartheta'\), that is \(v\vartheta'\ge3v-2\).
\item A pair with share \(\vartheta'\) has \(m'=(1-\vartheta')/2\). Under full orders at \(r\), the low-class sale
earns \(\rho\DT(r)m'-k\), because a sale never recruits. This is nonnegative only if
\(\vartheta'\le1-2k/(\rho\DT(r))=\bar\vartheta(r)\). At the benchmark, \(\bar\vartheta(r_1)=\reg{theta_max_full_eq_r1_B}\) and
\(\bar\vartheta(r_2)=\reg{theta_max_full_eq_r2_B}\). For example, the share \(0.9\) has \(\beta'=\reg{theta09_beta_B}\)
at the benchmark volume, but its sale at \(r_1\) earns \(\reg{theta09_lowsale_r1_B}\).
\item Strict recruitment at \(r_1\) and strict exclusion at \(r_2\) are \(\tau(r_1)<M'<\tau(r_2)\), that is
\(2\tau(r_1)-1<\vartheta'<2\tau(r_2)-1\), because \(\vartheta'=2M'-1\). At the benchmark volume
\(v=\reg{vol_full_B}\), \(v\le200/259\), so every \(u'\) in the band \((\reg{u_lo_frac},\reg{u_hi_frac})\) satisfies
\(3v\le2+u'(1-v)\), and part (a) gives a slot with \(\beta'\le1/3\), the same counts, and the same entry at \(r_1\) and at
\(r_2\). For every such slot, the weakest correct-order bound at \(r_1\) is
\(\rho\DT(r_1)[1-\tau(r_2)]-k=\reg{margin_obs_band_incentive_B}>0\), and \(\DT(r_2)>\DT(r_1)\); so full two-sided trade is
the unique outcome at both strengths.
\item In \(\mathcal F\), entry after the orders \(0\) and \(-1\) is \(\rho\), so \(\rho=\Pr(\text{preparation}\mid X\ne1)\).
Part (ii) of Proposition~\ref{prop:spread} gives \(I_-=\rho\DT\vartheta/2\), so \(\vartheta=2I_-/(\rho\DT(r))\). Then
\(v=\lambda+2b\) and \(\vartheta=\lambda/v\) give \(\lambda=v\vartheta\) and \(b=v(1-\vartheta)/2\), and
\(\beta=b/(1-\lambda)=v(1-\vartheta)/[2(1-v\vartheta)]\).\qedhere
\end{enumerate}
\end{proof}
The slot with \(u'=8/5\) has share \(\reg{alt_theta_B}\), lies in the band, and has the benchmark volume; its spread at
\(r_1\) is \(\reg{alt_ba_r1_B}\). The smallest share with the benchmark counts and \(\beta'\le1/3\) is
\(3-2/v\approx\reg{theta_min_counts_B}\). For example, the ratio \(u'=1/10\) and the benchmark volume give
\(\beta'>1/3\), so no slot of the model has this ratio and these counts.

At the share \(\bar\vartheta(r_1)=\reg{theta_max_full_eq_r1_B}\), the low-class sale earns exactly zero under full
orders. Rule IM admits the full sale, and full two-sided trade is then the unique outcome at \(r_1\). A sale rate below
\(\lambda'\) raises \(\mu(-1)\) above \(m'\) and makes the sale strictly profitable. A purchase rate below \(\lambda'\)
lowers \(\mu(1)\) below \(M'\) and makes the purchase strictly profitable. Rule I0 excludes the full sale
at this share, so this slot has no equilibrium at \(r_1\) under rule I0. Under rule IM, the shares that the benchmark counts
and the incentives at \(r_1\) allow therefore form the closed interval \([3-2/v,\bar\vartheta(r_1)]\). The values come
from \nolinkurl{P5_observables/p5_observables.py} (run log \nolinkurl{P5_observables/run_log.txt}).

\begin{corollary}[Counts under purchases only; analytical]\label{cor:purchasesonly}
Take the benchmark auction at the prior \(1/2\) and a slot in the band. In an equilibrium with \(\sigma_H=\delta_1\)
and \(\sigma_L=\delta_0\), the counts identify the slot: \(\lambda=2\E[X]\) and \(b=(v-\E[X])/2\), where
\(v=2b+\lambda/2\). Proposition~\ref{prop:coexistence} gives such an equilibrium exactly on \([r_F,r_N]\) (form F).
On branch R, class \(H\) buys with probability \(\sigma_H(1)<1\) and class \(L\) abstains. There the counts identify
the purchase rate \(s_+=\lambda\sigma_H(1)=2\E[X]\) and \(b\), but not \(\lambda\).
\end{corollary}

\begin{proof}[Proof of Corollary~\ref{cor:purchasesonly}]
At the prior \(1/2\), the order \(x\) has probability \([f_H(x)+f_L(x)]/2\). With \(\sigma_L=\delta_0\) and
\(s_+=\lambda\sigma_H(1)\), \(f_H(1)=b+s_+\), \(f_L(1)=b\) and \(f_H(-1)=f_L(-1)=b\). So \(\Pr(X=1)=b+s_+/2\),
\(\Pr(X=-1)=b\), \(\E[X]=s_+/2\) and \(v=2b+s_+/2\), which give \(s_+=2\E[X]\) and \(b=(v-\E[X])/2\). In form F,
\(s_+=\lambda\), which gives the first claim. By Proposition~\ref{prop:coexistence}, the high class buys with
probability one and the low class abstains exactly in form F on \([r_F,r_N]\).

On branch R,
\(s_+=b\,[\DT(r)/k-2]\) depends on the slot only through \(b\), and \(r_A\) does not depend on the slot. Take a slot
with the same \(b\) and a presence probability slightly below \(\lambda\). Its liquidity probability \(b/(1-\lambda)\)
falls, so it stays at most \(1/3\); its ratio stays in the band; and its end \(r_F\) stays above \(r\). So it has an R
equilibrium at \(r\) with the same counts, and the counts do not identify \(\lambda\) on branch R.
\end{proof}

Let \(A\) be the event that the challenger prepares and \(W\) the event that it acquires the target, and write
\(w_x=\Pr(W\mid X=x,A)\).

\begin{proposition}[Orders predict the winner; analytical]\label{prop:winner}
In the point-mass benchmark:
\begin{enumerate}[label=(\roman*),leftmargin=2em]
\item \(w_x=\mu(x)+[1-\mu(x)]\,\ell/r\), and \(\Pr(W\mid X=x)=e(x)\,w_x\).
\item Under full two-sided trade, \(w_1-w_{-1}=\vartheta(1-\ell/r)\), and this gap rises strictly with \(r\).
\item Under full two-sided trade, \(\E[X]=0\) and \(\Cov(X,\1_A)=(1-\rho)\chi_r(b+\lambda/2)\). For \(r>r_C\),
\(\Cov(X,\1_A)=0\) and \(\Cov(X,\1_W)=(b+\lambda/2)\,\rho\,\vartheta\,(1-\ell/r)>0\).
\end{enumerate}
\end{proposition}

\begin{proof}[Proof of Proposition~\ref{prop:winner}]
\begin{enumerate}[label=(\roman*),leftmargin=2em]
\item The cost is independent of the class and of \(R\), so preparation leaves the class posterior unchanged:
\(\Pr(H\mid X=x,A)=\mu(x)\). A prepared high-class challenger wins, because \(h>r\). A prepared low-class challenger
wins exactly when \(R<\ell\), with probability \(\ell/r\). Hence \(w_x=\mu(x)+[1-\mu(x)]\ell/r\), and
\(\Pr(W\mid X=x)=e(x)w_x\).
\item In \(\mathcal F\), \(w_1-w_{-1}=(M-m)(1-\ell/r)=\vartheta(1-\ell/r)\), and its derivative in \(r\) is
\(\vartheta\ell/r^2>0\). The function \(\mu\mapsto\mu+(1-\mu)\ell/r\) rises strictly, and the prices of \(\mathcal F\) rise
strictly with the posterior, so the win probability of a prepared challenger rises with the price.
\item In \(\mathcal F\), \(\E[X]=0\), so for any indicator \(Y\) with conditional means \(y_x=\E[Y\mid X=x]\),
\(\Cov(X,Y)=\E[XY]=\bar f(y_1-y_{-1})\). For \(Y=\1_A\), \(y_1-y_{-1}=e_+-\rho=(1-\rho)\chi_r\), so
\(\Cov(X,\1_A)=(1-\rho)\chi_r\bar f\). For \(Y=\1_W\), \(y_x=e(x)w_x\). For \(r>r_C\), \(e\equiv\rho\), so \(\Cov(X,\1_A)=0\) and
\(\Cov(X,\1_W)=\bar f\rho(w_1-w_{-1})=\bar f\rho\vartheta(1-\ell/r)>0\).\qedhere
\end{enumerate}
\end{proof}

Add a public event after the price that reveals only whether the challenger prepared, and nothing else. Market
makers then price again.

\begin{proposition}[Announcement returns; analytical]\label{prop:announcement}
Take any equilibrium at a strength \(r\).
\begin{enumerate}[label=(\roman*),leftmargin=2em]
\item After the order \(x\), the price changes by \(\Delta^1_x=[1-e(x)]\,A_r(\mu(x))\) when the event reports
preparation, and by \(\Delta^0_x=-e(x)\,A_r(\mu(x))\) when it reports no preparation. The second report has positive
probability exactly when \(e(x)<1\).
\item After an order \(x\) with \(e(x)<1\), the mean change over both reports, with their model probabilities, is zero.
\item Pool all orders, and take the sample of preparation reports. Its mean price change
\(\sum_x\Pr(X=x\mid A)\,\Delta^1_x\) is positive. After an order with \(e(x)=1\), the report of preparation changes
the price by zero.
\item Under full two-sided trade, the price after no order lies \((1-\rho)\chi_r(b+\lambda/2)A_r(M)\) below the expected
price before the order opportunity, although the posterior after no order is \(1/2\).
\end{enumerate}
\end{proposition}

\begin{proof}[Proof of Proposition~\ref{prop:announcement}]
After the order \(x\), the event reveals whether the challenger prepared.
\begin{enumerate}[label=(\roman*),leftmargin=2em]
\item The cost is independent of the class, so \(\Pr(H\mid x,\text{prepared})=\mu(x)\), and with preparation
\(\E[V_T]=t_j\) in class \(j\). So the price after a report of preparation is \(t_L+\DT\mu(x)=t_0+A_r(\mu(x))\). Without
preparation \(\E[V_T]=t_0\) in both classes. Subtract \(P(x)=t_0+e(x)A_r(\mu(x))\): \(\Delta^1_x=[1-e(x)]A_r(\mu(x))\) and
\(\Delta^0_x=-e(x)A_r(\mu(x))\). The report of no preparation has probability \(1-e(x)\), positive exactly when
\(e(x)<1\).
\item When \(e(x)<1\), \(e(x)\Delta^1_x+[1-e(x)]\Delta^0_x=e(x)[1-e(x)]A_r-[1-e(x)]e(x)A_r=0\).
\item Pooled over orders, a sample of preparation reports has mean change
\(\sum_x\Pr(X=x\mid A)\,[1-e(x)]A_r(\mu(x))\), with \(\Pr(X=x\mid A)=\Pr(X=x)e(x)/\Pr(A)\). Every term is nonnegative,
because \(A_r>0\) and \(e(x)\le1\); an order with \(e(x)=1\) contributes zero, since \(\Delta^1_x=0\). To find a positive
term, note that (E1) gives \(\sum_x\Pr(X=x)\mu(x)=1/2\), so some order \(x\) has \(\mu(x)\le1/2\). At the benchmark \(B_r(1/2)<c_H\) for \(1<r<h\), so \(\tau(r)>1/2\ge\mu(x)\), and
Lemma~\ref{lem:prices}(d) gives \(e(x)=\rho\) under both challenger rules. This order has
\(\Pr(X=x)\ge(1-\lambda)\eta(x)>0\), and its term \(\Pr(X=x)\rho(1-\rho)A_r(\mu(x))/\Pr(A)\) is positive. So the mean is
positive in every equilibrium; in \(\mathcal F\), the order \(-1\) is such an order.
\item In \(\mathcal F\), the expected price before the order opportunity is
\(\bar P=\sum_x\Pr(X=x)P(x)=t_0+\rho[\bar fA_r(m)+(1-2\bar f)A_r(1/2)+\bar fA_r(M)]+(1-\rho)\chi_r\bar fA_r(M)\). Since \(A_r\) is affine and
\(m+M=1\), \(\bar fA_r(m)+\bar fA_r(M)=2\bar fA_r(1/2)\), so \(\bar P=t_0+\rho A_r(1/2)+(1-\rho)\chi_r\bar fA_r(M)\). The price after no order
is \(t_0+\rho A_r(1/2)\), so \(P(0)-\bar P=-(1-\rho)\chi_r\bar fA_r(M)\), while \(\mu(0)=1/2\).\qedhere
\end{enumerate}
\end{proof}
In the benchmark at \(r_1\), the expected price is \(\reg{pbar_r1_B}\) and the price after no order is
\(\reg{p_neutral_r1_B}\). A sample of preparation reports has mean change \(\reg{evprep_mean_r1_B}\) at \(r_1\) and
\(\reg{evprep_mean_r2_B}\) at \(r_2\). The zero mean of part (ii) needs both reports with their model probabilities.

\paragraph{Two clusters at intermediate strengths.}
Let \(r_F<r<r_N\), so \(r<r_C\). Proposition~\ref{prop:coexistence} gives exactly two equilibria under rules IM and C1:
N, and F with \(\sigma_L=\delta_0\) (purchases only). In the second, \(\Pr(X=1)=b+\lambda/2\) and \(\Pr(X=-1)=b\), so
\(\E[X]=\lambda/2=\reg{imb_buy_B}\). The posterior after a sale is \(1/2\), and after no order it is
\(\mu_0=(1-2b-\lambda)/(2-4b-\lambda)\), with \(1/2-\mu_0=\lambda/[4(1-v_U)]>0\), \(v_U=2b+\lambda/2\). Both orders have
entry \(\rho\), so \(P(-1)-P(0)=\rho\DT\lambda/[4(1-v_U)]>0\); at \(r=\reg{str_r160}\) it is \(\reg{U_gap_r160_B}\). After a
purchase every challenger prepares. In N the three prices coincide and entry after a purchase is \(\rho\). Under rule
CM, the threshold pair T of Proposition~\ref{prop:coexistence} adds a third tuple: at \(r=\reg{str_r160}\) class \(H\)
buys with probability \(\reg{mix_sigma_r160_B}\), entry after a purchase is \(\reg{mix_e1_r160_B}\), the mean order is
\(\reg{mix_mean_order_r160_B}\), the low-class sale earns \(\reg{mix_lowsale_r160_B}\), and entry is
\(\reg{mix_entry_r160_B}\). The exhaustive two-cluster statement therefore uses rule C1.

\paragraph{Certificates.}
All values of this subsection come from the script
\nolinkurl{P5_observables/p5_observables.py}. It computes every value in
interval arithmetic at 60 digits from the exact input \(\beta=294655/10^6\) and an enclosure of \(\eu\); a sign counts only
when the enclosure excludes zero. The independent program
\nolinkurl{XP5/recompute.py} (run log
\nolinkurl{XP5/run_log.txt}) rebuilds every value from the auction cash flows
and the slot probabilities, with a rational Taylor bound on \(\eu\) through the term \(1/120!\), and it records one
comparison row per value in \nolinkurl{XP5/registry_recompute.csv}.
